\documentclass[11pt,a4paper]{article}
\usepackage{fontspec}
\setmainfont{Latin Modern Roman}
\usepackage[margin=23mm]{geometry}
\usepackage{amsmath,amssymb,amsthm,mathtools,mathrsfs}
\usepackage[unicode,hidelinks]{hyperref}
\hypersetup{pdftitle={The Original Observation and Its Mixed Factors},pdfauthor={Split-Zero research programme}}
\providecommand{\C}{\mathbb C}
\providecommand{\R}{\mathbb R}
\allowdisplaybreaks
\emergencystretch=2em
\title{The Original Observation and Its Mixed Factors\\[3pt]
\large Exact period vectors, primary coefficients and the retained kernel}
\author{Split-Zero research programme}
\date{14 September 2026}
\begin{document}
\maketitle
\paragraph{Portable source edition.}
Every TeX input used by this entry is expanded below. Historical file
and directory names in the preserved text are source locators; the
exact origin paths and hashes are recorded in \texttt{11\_SOURCE\_MAP.json}.
The other numbered proof files in this flat handoff supply the
accompanying complete calculations.\par

This separate source addendum continues the accepted Gamma mixed-row
calculation on the actual original observation. It retains the existing
source orders $1,k$, the original period matrix and cyclic projector,
all primary derivatives of the entire amplitude, and the full kernel.
The complete new calculation is followed by the exact proof dependencies
used here. Their full original editions also accompany this source in
\texttt{originals}; the selected proof spans retain their original bytes
and equation tags, as recorded in \texttt{SOURCE\_DEPENDENCIES.json}.

This editable entry is independent of the frozen main volume.
The original receiver interfaces are preserved in
\texttt{provenance/frozen\_receivers}; their next observation substitution
is proved below, with the arithmetic transition retained.
\tableofcontents
\clearpage

% BEGIN INLINE SOURCE 24; SHA256 d43c493eee0fc4e3e69f7383e7520b91900cdb183ed7c159bd9b2b473fdc5b04
% Exact original BRU observation, without a chosen replacement source order.
\section{The original period vectors and their Gamma generating map}
\label{sec:OPG}

The observation in this calculation is the literal BRU1--3 matrix. We
retain its entire amplitude unit, its primary multiplicities, its
period contours, its cyclic projector and its kernel. The only source
orders used are the already present orders $s\in\{1,k\}$. The letter
$s$ in a source-order subscript is an integer parameter; the coordinate
on the one-factor polynomial algebra will be denoted $z_1$ below to
keep it distinct. The generating variable is $z$.

\subsection{The original finite algebras and the observation}
Retain
\[
\begin{gathered}
 \rho=\tfrac12+\delta+i\gamma,\quad 0<\delta<\tfrac12,\quad\gamma>2,
 \quad\mathcal R=\{\rho,\bar\rho,1-\rho,1-\bar\rho\},\\
 h(z_1)=\prod_{\alpha\in\mathcal R}(z_1-\alpha)^m,
 \quad n=4m,\quad v_h=g/h,\quad g=2\xi,\quad m\geq1,\\
 E_h=\mathbb C[z_1]/(h),\qquad A_h=M_{z_1}:E_h\to E_h,
 \qquad U=M_{j_hv_h}:E_h\to E_h.
\end{gathered}\tag{OPG1}
\]
Here $m$ is the original complete zero order at each of the four
points. Consequently $v_h(\alpha)\ne0$ for every $\alpha\in\mathcal R$.
The entire-function remainder $j_hv_h$ records its derivatives of
orders $0,\ldots,m-1$ at every one of these points. In each primary
algebra it has nonzero constant term, and its inverse is the finite
geometric series for that constant plus a nilpotent. Thus $U$ is the
same invertible unit as in BRU, with none of its values prescribed.

For the original positive integer $k$, set
\[
\begin{gathered}
 c=k/2,\quad \ell_k=1+k(m-1),\quad q=(k+1)^2\ell_k,\\
 \chi(S)=\prod_{a,b=0}^k
   \bigl(S-c-(2a-k)\delta-i(2b-k)\gamma\bigr)^{\ell_k},\\
 E=\mathbb C[S]/(\chi),\quad A=M_S:E\to E,\quad
 Y=(A-cI)/i,\\
 A_\Sigma=\sum_{j=1}^k
      I_{E_h}^{\otimes(j-1)}\otimes A_h\otimes I_{E_h}^{\otimes(k-j)},
 \qquad
 \iota:E\hookrightarrow E_h^{\otimes k},\quad
 [P]\longmapsto[P(z_1+\cdots+z_k)].
\end{gathered}\tag{OPG2}
\]
For clarity, the stated map is injective with exactly this defining
polynomial. On an ordered primary tensor the eigenvalue of
$A_\Sigma$ is the sum of the ordered roots. Its nilpotent part is
the sum of the $k$ independent nilpotents of orders $m$.
Its power $k(m-1)$ has nonzero coefficient
$(k(m-1))!/((m-1)!)^k$ on their top monomial, and its next power
is zero. The distinct sums are exactly the $(k+1)^2$ displayed
roots of $\chi$, since the signs of the real and imaginary parts
can be selected independently in every factor. Distinct pairs
$(a,b)$ give distinct sums by $\delta,\gamma>0$. Their common
nilpotency index is $\ell_k$. These facts prove that the minimal
polynomial of $A_\Sigma$ is $\chi$, and therefore prove the
kernel of polynomial evaluation and the asserted injection.
In particular
\[
 \iota A=A_\Sigma\iota,\qquad
 \iota Y=\frac{A_\Sigma-cI}{i}\iota,\qquad
 \iota[1]=[1]^{\otimes k}.
 \tag{OPG3}
\]

Keep the original $u\ne0$, its logarithm branch, and $t=0$ in
the period construction. With $\Phi_h'=h$, $\Phi_h(0)=0$,
let $\ell_j$ be the ray of angle
$(\arg u+\pi+2\pi j)/(n+1)$ and let
$\Gamma_j=-\ell_0+\ell_j$, $1\leq j\leq n$, with the original
orientations. The original period map is
\[
 \Pi:E_h\to\mathbb C^n,\qquad
 (\Pi[P])_j=\int_{\Gamma_j}e^{\Phi_h(z_1)/u}
                    \operatorname{rem}_h P(z_1)\,dz_1.
 \tag{OPG4}
\]
The representative on the right has degree less than $n$. These
are exactly the original convergent polynomial period integrals;
their existence and full determinant are proved in PDE1--13 and
IPM1--7 in the retained source closure. This definition does not
assign a period to an arbitrary representative of a quotient class.
Let $C$ be the original matrix with rows corresponding to
$\Gamma_j\mapsto\Gamma_{j+1}-\Gamma_1$ for $j<n$ and
$\Gamma_n\mapsto-\Gamma_1$. The reduced endpoint basis gives
$C^{n+1}=I$. Hence the literal finite average below is idempotent:
\[
\begin{gathered}
 P_k=\frac1{n+1}\sum_{a=0}^{n}(C^{-a})^{\otimes k},\qquad
 L=P_k\Pi^{\otimes k}U^{\otimes k}\iota,\\
 B=\operatorname{im}L,\quad\jmath:B\hookrightarrow(\mathbb C^n)^{\otimes k},
 \qquad\Lambda:E\twoheadrightarrow B,\quad\jmath\Lambda=L,
 \qquad K=\ker\Lambda.
\end{gathered}\tag{OPG5}
\]
The basis inclusion $\jmath$, the inclusion $I:K\hookrightarrow E$,
and a coefficient section $S_*:B\to E$ with $\Lambda S_*=I_B$
are the fixed BRU choices. The full actual image is used at every
rank. Empty boundary matrices use the original empty conventions.

\subsection{An exact generating identity before and after observation}
Put $b_s=s/2$, $M_s=(2\pi)^{s/2}$, and retain the monic Gamma
polynomials and their full norms from GMB1:
\[
\begin{gathered}
 p_0(y)=1,\quad p_1(y)=y,\quad
 p_{d+1}(y)=yp_d(y)-d(b_s+d-1)p_{d-1}(y),\\
 h_d^{(s)}=M_s d!(b_s)_d,\quad
 r_d=p_d(Y)[1],\quad e_d=r_d/\sqrt{h_d^{(s)}},\quad
 \lambda_d=\Lambda e_d\in B.
\end{gathered}\tag{OPG6}
\]
The square root is the positive real square root. In particular
$e_d=[u_d]$ in the original notation, including the phase $i^{-d}$
in its leading coefficient. Define the holomorphic germs of
$\log(1+z^2)$ and $\arctan z$ by their values zero at $z=0$.
For $|z|<1$, the exact algebra-valued identity is
\[
 \sum_{d=0}^{\infty}r_d\frac{z^d}{d!}
 = (1+z^2)^{-s/4}\exp((\arctan z)Y)[1].
 \tag{OPG7}
\]
Indeed the right side is holomorphic in that disc, is $[1]$ at
zero, and satisfies
$(1+z^2)F'(z)=(Y-b_s zI)F(z)$.
Comparing its Taylor coefficients gives exactly the recurrence
in OPG6: the contribution of $z^2F'$ to coefficient $z^d/d!$
is $d(d-1)r_{d-1}$, and the contribution of $b_szF$ is
$b_sd r_{d-1}$. Thus its derivatives at zero are the $r_d$,
which proves the convergent identity and not merely a formal
coefficient equality.

All factors of $A_\Sigma$ act on different tensor coordinates
and therefore commute. Their matrix exponential power series
are absolutely convergent in finite dimension. The multinomial
expansion of every power consequently proves
\[
 \exp\left(w\frac{A_\Sigma-cI}{i}\right)[1]^{\otimes k}
 =\left(\exp\left(w\frac{A_h-\tfrac12I}{i}\right)[1]\right)^{\otimes k}.
 \tag{OPG8}
\]
The centre on the left is exactly $c=k/2$; each of the $k$
centres on the right is exactly $1/2$. Applying OPG3--5 to OPG7,
and then using OPG8, gives the original observed generating map
\[
\boxed{
 \jmath\sum_{d=0}^{\infty}\sqrt{h_d^{(s)}}\lambda_d\frac{z^d}{d!}
 = (1+z^2)^{-s/4}P_k
 \left[\Pi U\exp\left((\arctan z)
                \frac{A_h-\tfrac12I}{i}\right)[1]\right]^{\otimes k}.
}\tag{OPG9}
\]
The scalar outside $P_k$ has exponent $-s/4$. It has not acquired
an extra factor of $k$: OPG7 concerns one original Gamma source
on the sum coordinate, while OPG8 is the exact factorization of
that coordinate action in the already existing tensor algebra.
For $s=1$ and $s=k$ this gives the two stated source versions of
the same original observation. No additional order is selected.

\subsection{All primary coefficients and the original period integrals}
Here is a finite formula for every vector in OPG9. For
$\alpha\in\mathcal R$ write
\[
\begin{gathered}
 h_\alpha(z_1)=h(z_1)/(z_1-\alpha)^m,\qquad
 t_\alpha(z_1)=\sum_{j=0}^{m-1}
       \frac{(h_\alpha^{-1})^{(j)}(\alpha)}{j!}(z_1-\alpha)^j,\\
 \varepsilon_\alpha=[h_\alpha t_\alpha]\in E_h,
 \qquad E_{\alpha,a}(z_1)=
   \operatorname{rem}_h\bigl(h_\alpha(z_1)t_\alpha(z_1)
                                  (z_1-\alpha)^a\bigr),\quad0\leq a<m,\\
 \mathbf b_{\alpha,a}=\Pi[E_{\alpha,a}],\qquad
 (\mathbf b_{\alpha,a})_j=
   \int_{\Gamma_j}e^{\Phi_h(z_1)/u}E_{\alpha,a}(z_1)\,dz_1.
\end{gathered}\tag{OPG10}
\]
The Taylor inverse exists because $h_\alpha(\alpha)\ne0$.
Reduction modulo $(z_1-\alpha)^m$ gives
$\varepsilon_\alpha=1$ on its primary factor and zero on every
other primary factor. These observations prove
$\varepsilon_\alpha\varepsilon_\beta=0$ for $\alpha\ne\beta$,
$\varepsilon_\alpha^2=\varepsilon_\alpha$, and
$\sum_\alpha\varepsilon_\alpha=1$: their differences are divisible
by all four coprime factors and hence by $h$. They also prove that
$[E_{\alpha,a}]=\varepsilon_\alpha(z_1-\alpha)^a$ is a basis of
$E_h$. Every representative in OPG10 has degree less than $n$,
so each integral is precisely of the original polynomial type.

For $w\in\mathbb C$ put
$\mathbf v(w)=\Pi U\exp(w(A_h-\tfrac12I)/i)[1]$.
The primary nilpotent $N_\alpha=(z_1-\alpha)\varepsilon_\alpha$
satisfies $N_\alpha^m=0$. Taylor multiplication in that algebra
therefore proves, with all amplitude derivatives retained,
\[
\boxed{
 \mathbf v(w)=\sum_{\alpha\in\mathcal R}
 e^{w(\alpha-1/2)/i}\sum_{a=0}^{m-1}\mathbf b_{\alpha,a}
       \sum_{\nu=0}^{a}
       \frac{v_h^{(\nu)}(\alpha)}{\nu!}
       \frac{(w/i)^{a-\nu}}{(a-\nu)!}.
}\tag{OPG11}
\]
For a direct verification before integration, both
$U\exp(w(A_h-\tfrac12I)/i)[1]$ and
$j_h(v_h(z_1)e^{w(z_1-1/2)/i})$ have the same first $m$ Taylor
coefficients at every $\alpha$. Their difference is therefore
zero in $E_h$. The representative to which $\Pi$ is applied is
exactly the finite polynomial in the basis $E_{\alpha,a}$ whose
coefficients are displayed in OPG11. In particular this proof
does not replace the integrand in OPG10 by the unreduced entire
function $v_h(z_1)e^{w(z_1-1/2)/i}$. The precise map from that
entire function to the period vector is its finite Taylor
remainder $j_h$, followed by the degree-less-than-$n$ representative,
followed by the original polynomial period integral.

The cyclic projector is also fully evaluated as a finite operation:
\[
 \jmath\sum_{d\geq0}\sqrt{h_d^{(s)}}\lambda_d\frac{z^d}{d!}
 =\frac{(1+z^2)^{-s/4}}{n+1}
       \sum_{a=0}^{n}\bigl(C^{-a}\mathbf v(\arctan z)\bigr)^{\otimes k}.
 \tag{OPG12}
\]
This is obtained by applying each tensor power in OPG5 to a pure
tensor. It retains every term of the original average, its factor
$1/(n+1)$, and its inverse powers of $C$.

There is also a finite coefficient formula requiring no series
extraction. For an ordered $k$-tuple
$\boldsymbol\alpha\in\mathcal R^k$ set
$\omega_{\boldsymbol\alpha}=(\sum_i\alpha_i-c)/i$.
For $0\leq a_i<m$ and $0\leq\nu_i\leq a_i$ put
$D(\mathbf a,\boldsymbol\nu)=\sum_i(a_i-\nu_i)$ and
\[
 T_d(\boldsymbol\alpha,\mathbf a)=
 \sum_{\substack{0\leq\nu_i\leq a_i\\1\leq i\leq k}}
 \frac{i^{-D(\mathbf a,\boldsymbol\nu)}
        p_d^{(D(\mathbf a,\boldsymbol\nu))}(\omega_{\boldsymbol\alpha})}
        {\prod_{i=1}^k(a_i-\nu_i)!}
 \prod_{i=1}^k\frac{v_h^{(\nu_i)}(\alpha_i)}{\nu_i!}.
 \tag{OPG13}
\]
As usual the polynomial derivative is zero when its order exceeds
$d$. Multivariate Taylor expansion in the independent nilpotents
$(z_i-\alpha_i)$ proves the exact formula
\[
\boxed{
 \jmath\lambda_d=
 \frac1{\sqrt{h_d^{(s)}}}\frac1{n+1}
 \sum_{a=0}^{n}\ \sum_{\boldsymbol\alpha\in\mathcal R^k}
 \sum_{\substack{0\leq a_i<m\\1\leq i\leq k}}
 T_d(\boldsymbol\alpha,\mathbf a)
 \bigotimes_{i=1}^k C^{-a}\mathbf b_{\alpha_i,a_i}.
}\tag{OPG14}
\]
Indeed on that primary tensor, expand each $v_h(z_i)$ first.
For the remaining multi-degree $a_i-\nu_i$, Taylor expansion of
$p_d((\sum_i z_i-c)/i)$ supplies its derivative of the total
order $D$, the factor $i^{-D}$, and the factorials
$\prod_i(a_i-\nu_i)!$. These are exactly the coefficients in
OPG13. Apply the $k$ polynomial period maps and the literal cyclic
average to get OPG14. All sums are finite. Ordered tuples have
not been replaced by just their eigenvalue sum, so both the
amplitude data and every independent primary nilpotent survive.

\subsection{The actual observation increment without a coordinate inverse}
Retain the GMB maps
$C_s=[e_0,\ldots,e_{q-1}]:\mathbb C^q\to E$,
$v_j=C_s^{-1}e_{q+j}$,
$K_j=I_q+\sum_{a<j}v_av_a^*$, and $R_s=\Lambda C_s$.
The low-degree leading coefficients prove $C_s$ invertible as in
GMB2. On each column, applying the maps gives
\[
 R_s v_j=\lambda_{q+j},\qquad
 R_sK_jR_s^*=\sum_{d=0}^{q+j-1}\lambda_d\lambda_d^*.
 \tag{OPG15}
\]
For the second equality the identity term in $K_j$ gives exactly
the first $q$ columns of $R_s$, and each remaining rank-one term
gives the next displayed column. The low vectors $\lambda_d$,
$0\leq d<q$, span $B$ because $C_s$ is invertible and $\Lambda$
is onto. Thus, in the fixed original basis of $B$, the matrix
\[
 \mathcal M_D=\sum_{d=0}^{D-1}\lambda_d\lambda_d^*\quad(D\geq q)
 \tag{OPG16}
\]
is positive definite whenever $B$ has positive dimension:
for a nonzero covector its quadratic form is the sum of its
squared values on spanning vectors. Consequently the exact
original boundary increment of GMB8 is
\[
 \boxed{\displaystyle
 \xi_j=\lambda_{q+j}^*\mathcal M_{q+j}^{-1}\lambda_{q+j},\qquad
 1+\xi_j=\frac{\det\mathcal M_{q+j+1}}{\det\mathcal M_{q+j}}.}
 \tag{OPG17}
\]
The determinant equality follows by conjugating the rank-one
update by $\mathcal M_{q+j}^{-1/2}$; its sole possible nonzero
added eigenvalue is the displayed quadratic form. If $B=0$,
the vector is empty, $\xi_j=0$ and both determinants equal one.
Thus OPG11--14 compute the boundary increment from the original
period vectors, with no $C_s^{-1}$ in OPG17. This cancellation
is an identity of the proved maps, rather than a change of the
observation or the source measure.

At the four original endpoints the boundary contribution is
therefore the exact finite expression
\[
 F_B^{(s)}=
 \log\frac{\det\mathcal M_{2q}\det\mathcal M_{2q+1}}
                 {\det\mathcal M_q\det\mathcal M_{q+1}}
 =\log(1+\xi_0)+2\sum_{j=1}^{q-1}\log(1+\xi_j)
                         +\log(1+\xi_q).
 \tag{OPG18}
\]
Expanding the last sum of consecutive log determinant differences
leaves coefficient $-1$ at $\mathcal M_q,\mathcal M_{q+1}$ and
$+1$ at $\mathcal M_{2q},\mathcal M_{2q+1}$, proving every sign.
GMB13 proves that the other original factor is
$(1+t_j)/(1+\xi_j)$, with $t_j=v_j^*K_j^{-1}v_j$ and
$0\leq\xi_j\leq t_j$. Thus OPG18 is the actual BRU boundary
contribution, and the full original kernel contribution remains
$\mathcal W_q(\log((1+t_j)/(1+\xi_j)))$.

\subsection{A finite recurrence retaining the exact action defect}
The period formula also has a finite recurrence on the fixed
original coefficient spaces. Define the uniquely typed projection
$\kappa:E\to K$ by
$I\kappa=I_E-S_*\Lambda$. It exists because the right side has
image in $\ker\Lambda=\operatorname{im}I$, and $I$ is an
inclusion. It satisfies $\kappa I=I_K$ and $\kappa S_*=0$.
In particular $E=I(K)\oplus S_*(B)$ as coefficient spaces, without
an assertion that this splitting is preserved by $Y$.
Set
\[
\begin{gathered}
 \mathsf Y_B=\Lambda YS_*:B\to B,\quad
 \mathsf D=\Lambda YI:K\to B,\\
 \mathsf C=\kappa YS_*:B\to K,\quad
 \mathsf Y_K=\kappa YI:K\to K,\quad
 \zeta_d=\kappa e_d\in K.
\end{gathered}\tag{OPG19}
\]
All four maps are computable from OPG1--5 and the fixed BRU
section. For $d\geq0$ put
\[
 a_d=\frac1{\sqrt{(d+1)(b_s+d)}},\qquad
 b_d=\sqrt{\frac{d(b_s+d-1)}{(d+1)(b_s+d)}}\quad(d\geq1),
 \qquad b_0=0.
 \tag{OPG20}
\]
Set $\lambda_{-1}=0$ and $\zeta_{-1}=0$. Then the exact coupled
recurrence is
\[
\boxed{\begin{aligned}
 \lambda_{d+1}&=a_d(\mathsf Y_B\lambda_d+\mathsf D\zeta_d)
                                          -b_d\lambda_{d-1},\\
 \zeta_{d+1}&=a_d(\mathsf C\lambda_d+\mathsf Y_K\zeta_d)
                                          -b_d\zeta_{d-1},\\
 \lambda_0&=\Lambda[1]/\sqrt{M_s},\qquad
 \zeta_0=\kappa[1]/\sqrt{M_s}.
\end{aligned}}\tag{OPG21}
\]
To prove this, the norm ratio from OPG6 gives
$\sqrt{h_d^{(s)}/h_{d+1}^{(s)}}=a_d$ and
$d(b_s+d-1)\sqrt{h_{d-1}^{(s)}/h_{d+1}^{(s)}}=b_d$.
Hence $e_{d+1}=a_dYe_d-b_de_{d-1}$, including $d=0$ by its
initial values. Substituting $e_d=S_*\lambda_d+I\zeta_d$ and
applying $\Lambda$ and $\kappa$ proves both lines of OPG21.
It retains the kernel coordinate at every degree and supplies
the period vectors through degree $2q$, exactly the range used
in OPG17--18.

The failure of a putative boundary action has the precise map
\[
 \Lambda Y-\mathsf Y_B\Lambda=\mathsf D\kappa:E\to B,\qquad
 \jmath\mathsf D=
 \frac1iP_k\Pi^{\otimes k}U^{\otimes k}
                   (A_\Sigma-cI)\iota I:K\to(\mathbb C^n)^{\otimes k}.
 \tag{OPG22}
\]
The first equality follows from $I_E=S_*\Lambda+I\kappa$;
the second follows from OPG3--5. These equations preserve the
cyclic average and the full unit. They also prove exactly that
$Y(K)\subseteq K$ holds if and only if $\mathsf D=0$.
No such vanishing is used anywhere in OPG7--21. Thus every
possible boundary action defect is retained inside a proved
coupled action and inside the explicit original period formula.

For completeness the multiplication action does descend to a
canonical quotient that remembers every observed iterate. Define
\[
 \mathcal K_\infty=\bigcap_{a=0}^{q-1}\ker(\Lambda Y^a),\qquad
 \mathcal O:E\to B^{\oplus q},\quad
 x\longmapsto(\Lambda x,\Lambda Yx,\ldots,\Lambda Y^{q-1}x).
 \tag{OPG23}
\]
The polynomial $\chi(c+iT)$ annihilates $Y$ by OPG2, has degree
$q$ and leading coefficient $i^q\ne0$. It therefore expresses
$Y^q$ as a linear combination of $I,Y,\ldots,Y^{q-1}$, with
the exact coefficients of this original polynomial. If
$x\in\mathcal K_\infty$, each component of $\mathcal O(Yx)$
vanishes by this relation. Thus $\mathcal K_\infty$ is
$Y$-invariant and is precisely $\bigcap_{a\geq0}\ker(\Lambda Y^a)$.
Moreover any $Y$-invariant subspace contained in $K$ lies in
this intersection. Consequently the first isomorphism map
\[
\begin{gathered}
 E/\mathcal K_\infty\ \xrightarrow{\ \overline{\mathcal O}\ }
 \operatorname{im}\mathcal O,\quad[x]\longmapsto\mathcal O(x),\\
 \chi_j=[T^j]\chi(c+iT),\quad\chi_q=i^q,\\
 \mathsf{Shift}(b_0,\ldots,b_{q-1})=
 \left(b_1,\ldots,b_{q-1},-i^{-q}\sum_{j=0}^{q-1}\chi_j b_j\right).
\end{gathered}
 \tag{OPG24}
\]
is a bijection intertwining multiplication by $Y$ with the shift
of the $q$ components, whose last component is the stated exact
linear combination. It projects to the original observation by
the first component, with retained kernel $K/\mathcal K_\infty$.
This proves an action-preserving enlargement of the observed
data on the original algebra. It neither discards the original
kernel in OPG21 nor presumes a vanished defect in OPG22.

The original period isomorphism also locates that defect in the
cyclic projection. Its inverse exists by the retained period
determinant theorem used in OPG4. On the unchanged ambient tensor
space put
\[
 \mathcal V=\Pi^{\otimes k}U^{\otimes k}\iota,
 \qquad
 \mathcal T=\sum_{j=1}^k
 \left(\Pi\frac{A_h-\tfrac12I}{i}\Pi^{-1}\right)_j.
 \tag{OPG25}
\]
The subscript $j$ means that the displayed factor acts in position
$j$ and identities act in all other positions. Multiplication by
$j_hv_h$ commutes with $A_h$. Equations OPG3 and OPG25 therefore
give $\mathcal VY=\mathcal T\mathcal V$ exactly. Since
$L=P_k\mathcal V$ and $P_k^2=P_k$, splitting the ambient identity
as $P_k+(I-P_k)$ now proves
\[
\begin{gathered}
 LY=P_k\mathcal T L+P_k\mathcal T(I-P_k)\mathcal V,\\
 \jmath\mathsf D=P_k\mathcal T(I-P_k)\mathcal V I.
\end{gathered}\tag{OPG26}
\]
For the second equality, $P_k\mathcal V I=LI=0$, so the
first term of the first equality vanishes on $I(K)$. This is
the same defect as OPG22 on its original domain. The two ambient
terms in the first line sum to a vector in $\jmath(B)$; neither
individual term is asserted to preserve that image. OPG19--21
give the full recurrence within $K\oplus B$ without such an
assertion. Thus the preserved unit and period isomorphism carry
the multiplication action exactly, and OPG26 identifies precisely
the retained contribution of the original cyclic projection.

Equations OPG9, OPG11--14, OPG17 and OPG21 now give the same
original boundary increments by analytic generating functions,
finite primary coefficients, and finite coupled recurrences.
Their derivations establish the exact maps connecting these
calculations. In particular the new input into the original
mixed-control problem is the fully specified sequence
$\lambda_0,\ldots,\lambda_{2q}$ with its actual unit and periods,
rather than a freely selected observation or a source-order
replacement. The formulas themselves do not assign a growing-
parameter estimate to the individual weighted sums; they supply
their unchanged mathematical input for that calculation.

% END INLINE SOURCE 24

\clearpage

% BEGIN INLINE SOURCE 25; SHA256 b7de12c81133c7fb75e1bd3e46be3ca591988692ff6f6b8ef6ddecfbd4bfa92e
\section{The exact receiving map into the original mixed factors}
\label{sec:OPR}

The current receiver MRI1--5 is retained as a frozen predecessor.
This section proves its next observation input from OPG1--26.
Fix either original source order $s\in\{1,k\}$, put $D=q+j$,
and retain all the fixed coefficient maps of OPG and GMB.
In particular $e_D=[u_D]$, $\lambda_D=\Lambda e_D$, and
$\mathcal M_D=\sum_{d<D}\lambda_d\lambda_d^*$ on the original
boundary image. At the preceding polynomial degree $D-1$, GMB4
and OPG15 give
\[
\begin{gathered}
 G_{D-1}^{(s)}=C_s^{-*}K_j^{-1}C_s^{-1},\quad
 (G_{D-1}^{(s)})^{-1}=C_sK_jC_s^*,\\
 Q_{D-1}^{(s)}=\mathcal M_D^{-1},\qquad
 S_{D-1}^{(s)}=C_sK_jC_s^*\Lambda^*\mathcal M_D^{-1}:B\to E.
\end{gathered}\tag{OPR1}
\]
Indeed $\Lambda(G_{D-1}^{(s)})^{-1}\Lambda^*=
R_sK_jR_s^*=\mathcal M_D$. Direct multiplication therefore
proves $\Lambda S_{D-1}^{(s)}=I_B$ and
$I^*G_{D-1}^{(s)}S_{D-1}^{(s)}=0$. Thus this is precisely
the original minimum section for this degree and this source.
All statements use their unique empty meaning if $B=0$.

MRI2 uses the kernel vector
$r_j^{\rm ker}=v_j-K_jR_s^*\mathcal M_D^{-1}\lambda_D$
in the $C_s$ coordinates. The fixed-section kernel vector of
OPG21 is $\zeta_D=\kappa e_D$, with $I\kappa=I_E-S_*\Lambda$.
The exact morphism between these two specified vectors is
\[
\boxed{\begin{aligned}
 C_s r_j^{\rm ker}
   &=e_D-S_{D-1}^{(s)}\lambda_D\\
   &=I\zeta_D+(S_*-S_{D-1}^{(s)})\lambda_D
    =I\vartheta_{D,j}^{(s)},\\
 \vartheta_{D,j}^{(s)}
   &=\zeta_D-\kappa S_{D-1}^{(s)}\lambda_D\in K.
\end{aligned}}\tag{OPR2}
\]
For the first equality substitute $v_j=C_s^{-1}e_D$ and
$R_s^*=C_s^*\Lambda^*$ into MRI2. For the second substitute
$e_D=I\zeta_D+S_*\lambda_D$. The difference of the two
sections lands in $K$ because their observations are both $I_B$.
Applying $\kappa$ and using $\kappa S_*=0$ proves the last
line and hence the third equality. This proves the complete
section correction in the original coefficient space.

Retain the actual kernel Gram
$H_{D-1}^{(s),K}=I^*G_{D-1}^{(s)}I$. Congruencing OPR1 by
$C_s$ and then using OPR2 proves
\[
 t_j-\xi_j
 = (r_j^{\rm ker})^*K_j^{-1}r_j^{\rm ker}
 = (\vartheta_{D,j}^{(s)})^*
                    H_{D-1}^{(s),K}\vartheta_{D,j}^{(s)}.
 \tag{OPR3}
\]
The first equality is the proved MRI2/GMB10 orthogonal
decomposition. Thus the recurrence OPG21 supplies both the
actual boundary column and, through the complete section
correction OPR2, the exact kernel norm needed at each step.
The original ordered factors are consequently
\[
\begin{gathered}
 \xi_j=\lambda_D^*\mathcal M_D^{-1}\lambda_D,\qquad
 1+\frac{\nu_{D-1}}{a_{D-1}+\beta_{D-1}}=1+\xi_j,\\
 1+\frac{\beta_{D-1}}{a_{D-1}}
 =\frac{1+t_j}{1+\xi_j}
 =1+\frac{(\vartheta_{D,j}^{(s)})^*
                  H_{D-1}^{(s),K}\vartheta_{D,j}^{(s)}}{1+\xi_j}.
\end{gathered}\tag{OPR4}
\]
The first line uses OPG17 and GMB13. The second line follows
by inserting OPR3 into that same GMB13 identity. This preserves
the denominator $1+\xi_j$ and the original order of elimination.

The finite endpoint operation remains
$\mathcal W_q(f)=f_0+2\sum_{j=1}^{q-1}f_j+f_q$.
Applying it to the logarithms in OPR4 reproduces MRI4 on its
unchanged source. The exact arithmetic receiving equality is
therefore still
\[
\begin{aligned}
 \mathcal B_k^{\rm ar}
 &=F_K^{(k)}+F_B^{(k)}+\mathcal H_k+\delta_k^\sigma\\
 &=F_K^{(1)}+F_B^{(1)}+\delta_k^\sigma.
\end{aligned}\tag{OPR5}
\]
Its proof is substitution into the existing MRI5/BRI2 source
identity: $F_K^{(k)}+F_B^{(k)}=\mathcal B_k^0$,
$F_K^{(1)}+F_B^{(1)}=\mathcal B_k^\sigma$,
$\mathcal H_k=\mathcal B_k^\sigma-\mathcal B_k^0$, and
$\delta_k^\sigma=\mathcal B_k^{\rm ar}-\mathcal B_k^\sigma$.
No term is assigned a new sign or asymptotic coefficient.
The combined coefficient received by BRI6 and BRI9 remains
attached to this same sum. OPR1--4 provide its individual
original factors from the actual OPG observation vectors,
with all original source transitions present in OPR5.

% END INLINE SOURCE 25

\clearpage
\appendix

% BEGIN INLINE SOURCE 26; SHA256 0917a0a4c54aeb40e694240dd697c712c01b8b4aaabd53cf00ae3f51da16a311
% Canonical intrinsic operation: no independent source order is inserted.
\section{The Gamma orders produced by the existing exterior and sum maps}
\label{sec:IGO}

The source order in this section is one of the two orders already present
in the construction. The determinant size is the actual exterior degree.
Every relation coefficient remains in the resulting density.

\subsection{Original coordinates and the exact source measure}
Retain
\[
 k=4l+1,\quad l,m\geq1,\quad e=1+k(m-1),\quad
 q=e(k+1)^2,\quad c=k/2,\quad
 0<\delta<1/2,\quad\gamma>2,
\]
\[
 \chi(S)=\prod_{a,d=0}^{k}
 [S-c-(2a-k)\delta-i(2d-k)\gamma]^e,
 \qquad s\in\{1,k\},\quad b=s/2,\quad M_s=(2\pi)^{s/2}.
 \tag{IGO1}
\]
Here $s=1$ is the fixed source $\sigma$ and $s=k$ is the original
tensor Gamma source. The original line throughout is $S=c+iy$.
The change of letter from the second root index to $d$ leaves every
root and its order unchanged.

For each real $v>0$ define a density on the real coordinate by
\[
 \rho_v(y)=\frac{2^v}{4\pi\Gamma(v)}
       \left|\Gamma(v/2+iy/2)\right|^2,
 \qquad dm_{0,s}(y)=M_s\rho_b(y)\,dy.
 \tag{IGO2}
\]
Thus the full mass $M_s$ is retained. In particular
$dm_{0,1}=|\Gamma(1/4+iy/2)|^2dy/(2\pi)=d\sigma$.
For the other existing source, the recurrence of the Gamma function
at $k/4=l+1/4$ gives
\[
 dm_{0,k}(y)=\beta_k\prod_{j=0}^{l-1}
           [y^2+(2j+\tfrac12)^2]d\sigma(y),\qquad
 \beta_k=\frac{(2\pi)^{k/2}}{\sqrt2\Gamma(k/2)}.
\]
Indeed the squared Gamma recurrence contributes $4^{-l}$ times
the displayed product. The coefficient in IGO2, multiplied by
$4^{-l}$ and divided by the coefficient $1/(2\pi)$ of $\sigma$,
is $M_k2^{k/2-1-2l}/\Gamma(k/2)=\beta_k$.
This proves agreement with the original tensor source including
its full polynomial factor and mass.
For real $|t|<\pi/2$ and real $\xi$, the exact transforms are
\[
 \int_{\mathbb R}e^{ty}\rho_v(y)\,dy=(\cos t)^{-v},\qquad
 \int_{\mathbb R}e^{i\xi y}\rho_v(y)\,dy=(\cosh\xi)^{-v},
 \qquad \phi_s(t):=\int e^{ty}dm_{0,s}=M_s(\cos t)^{-b}.
 \tag{IGO3}
\]

Here is a direct justification of the source identities and the
integrability used below. The beta integral with the substitution
$u=\exp(2x)$ gives
\[
 \frac{\Gamma(v/2+iy/2)\Gamma(v/2-iy/2)}{\Gamma(v)}
  =2\int_{\mathbb R} e^{iyx}(2\cosh x)^{-v}\,dx.
\]
Consequently $\rho_v(y)=(2\pi)^{-1}
\int e^{iyx}(\cosh x)^{-v}dx$. Fourier inversion applies: the
function and its first two derivatives are integrable, and two
integrations by parts make its transform integrable. This proves
the characteristic identity and mass one in IGO3, with the displayed
Fourier convention. More strongly, for every $0<\theta<\pi/2$, move
the $x$ contour to $x+i\theta\operatorname{sgn}(y)$. The analytic
branch of $(\cosh x)^{-v}$ is fixed by its positive values on the
real axis; it exists in this strip because $\cosh x$ has no zero
there. Its modulus decays as $\exp(-v|\operatorname{Re}x|)$ on
each closed smaller strip. The vertical sides of the rectangle
therefore tend to zero and
$\rho_v(y)\leq C_{v,\theta}\exp(-\theta|y|)$.
This proves absolute convergence of every polynomial moment with
an exponential factor $e^{ty}$ when $|t|<\pi/2$. Holomorphic
continuation of the characteristic identity through this strip
then gives the first identity in IGO3. Strict positivity follows
from the zero-free Euler Gamma product in $\operatorname{Re}z>0$:
\[
 \frac1{\Gamma(z)}=z\exp(\gamma_{\!E}z)
       \prod_{n=1}^{\infty}(1+z/n)\exp(-z/n).
\]
Indeed the logarithms of the factors after their linear terms are
removed converge locally uniformly, and none of the factors vanishes
there. This product follows by taking the beta-integral limit
$\Gamma(z)=\lim_{n\to\infty}n!n^z/[z(z+1)\cdots(z+n)]$;
the linear sum is absorbed by
$\gamma_{\!E}=\lim_n(\sum_{j=1}^n1/j-\log n)$.


% END INLINE SOURCE 26


% BEGIN INLINE SOURCE 27; SHA256 8d5957fb08f2e1d6448a8130dc4cc5ff530cfae2791ef15d24fd274fceaf4542
\section{The full norms of the two existing Gamma sources}
\label{sec:OGN}

The included canonical source proof IGO1--3 defines
\[
 dm_{0,s}(y)=M_s\frac{2^{b_s}}{4\pi\Gamma(b_s)}
       |\Gamma(b_s/2+iy/2)|^2\,dy,
 \qquad s\in\{1,k\},\quad b_s=s/2,\quad M_s=(2\pi)^{s/2},
 \tag{OGN1}
\]
and proves its exact Laplace transform
$\int e^{ty}dm_{0,s}=M_s(\cos t)^{-b_s}$ for real $|t|<\pi/2$.
Its proof includes exponential domination on each smaller strip,
so all the differentiated integrals below converge absolutely.
Every source mass and the original line $S=c+iy$, $c=k/2$, are
retained here.

Let $p_d$ be the recurrence polynomials of OPG6, and put
$G(z,y)=(1+z^2)^{-b_s/2}e^{y\arctan z}$ near zero. The
coefficient calculation in OPG7, with scalar $y$ in place of its
matrix, proves $G(z,y)=\sum_d p_d(y)z^d/d!$.
Each $p_d$ is real and monic by the recurrence and its initial
values. For sufficiently small real $z,w$, the elementary identity
\[
 \cos(\arctan z+\arctan w)
       =\frac{1-zw}{\sqrt{(1+z^2)(1+w^2)}}
\]
and the full Laplace transform give
\[
 \int_{\mathbb R}G(z,y)G(w,y)\,dm_{0,s}(y)
                  =M_s(1-zw)^{-b_s}.
 \tag{OGN2}
\]
Choose the neighbourhood so that
$|\arctan z|+|\arctan w|\leq a<\pi/2$.
Every fixed mixed derivative of the integrand is bounded on a
smaller neighbourhood by a constant times a polynomial in $|y|$
times $e^{a|y|}$. The exponential domination proved in IGO3,
with a strictly larger exponent smaller than $\pi/2$, integrates
this bound. Differentiation under the integral is therefore valid
at every fixed order. Applying $\partial_z^d\partial_w^e$ at zero
to OGN2 yields
\[
 \int p_d(y)p_e(y)\,dm_{0,s}(y)
  =\begin{cases}0,&d\ne e,\\
        M_s d!(b_s)_d,&d=e.
    \end{cases}
 \tag{OGN3}
\]
Indeed $(1-zw)^{-b_s}=\sum_{a\geq0}(b_s)_a(zw)^a/a!$,
so its $(d,e)$ derivative at zero vanishes off the diagonal and
equals $d!(b_s)_d$ on it. Since $b_s>0$, every diagonal entry
is strictly positive. Consequently
$u_d(S)=p_d((S-c)/i)/\sqrt{M_s d!(b_s)_d}$ is exactly the
orthonormal source polynomial used in GMB1 and OPG6, with
leading coefficient $i^{-d}/\sqrt{M_s d!(b_s)_d}$ in $S$.
This proves those source norms directly from the two retained
source measures and fixes the full phase and mass of each
observation column.

% END INLINE SOURCE 27

\clearpage

% BEGIN INLINE SOURCE 28; SHA256 32c56427cb2f62c0d41bd2f6830813a84aa01b1bff21a39fd4291e9b54f35556
% Original source remainder coordinates and the full BRU mixed row.
\section{The exact map from Gamma remainders to the original mixed rows}
\label{sec:GMB}

This calculation uses only the two sources already present in the
original construction. It retains the packet order $k$, its centre
$c=k/2$, its monic polynomial $\chi$ of degree $q$, and the original
surjection $\Lambda:E_\chi\to B_{\rm obs}$, where
$E_\chi=\mathbb C[S]/(\chi)$. The source order is $s\in\{1,k\}$.
Write $b_s=s/2$ and $M_s=(2\pi)^{s/2}$. For the full source measure
$dm_{0,s}$, the monic real orthogonal polynomials $p_d(y)$ and their
squared norms are
\[
\begin{gathered}
 p_0=1,\quad p_1=y,\qquad
 p_{d+1}(y)=yp_d(y)-d(b_s+d-1)p_{d-1}(y),\\
 h_d^{(s)}=M_s d!(b_s)_d,\qquad
 u_d(S)=\frac{p_d((S-c)/i)}{\sqrt{h_d^{(s)}}}.
\end{gathered}\tag{GMB1}
\]
The recurrence and every source mass are those proved for the
original Gamma source. On $S=c+iy$ the $u_d$ are orthonormal;
their leading coefficients in $S$ are $i^{-d}/\sqrt{h_d^{(s)}}$.
Thus no phase or mass is removed by introducing these explicit
coordinates. Inner products below are conjugate linear in the
first variable.

Let $C_s:\mathbb C^q\to E_\chi$ have columns
$[u_0],\ldots,[u_{q-1}]$ in the fixed original remainder frame.
It is invertible, with the literal determinant
\[
 \det C_s=i^{-q(q-1)/2}\prod_{d=0}^{q-1}(h_d^{(s)})^{-1/2}.
 \tag{GMB2}
\]
For $j\geq0$ define the actual remainder coordinates and matrices
\[
\begin{gathered}
 v_j=C_s^{-1}[u_{q+j}],\qquad
 K_j=I_q+\sum_{a=0}^{j-1}v_av_a^*,\qquad
 t_j=v_j^*K_j^{-1}v_j,\qquad d_j^{\rm gain}=1+t_j,\\
 A_j=[I_q,v_0,\ldots,v_{j-1}]:\mathbb C^{q+j}\to\mathbb C^q.
\end{gathered}\tag{GMB3}
\]
The empty sum is zero. In particular $K_0=I_q$, $K_j=A_jA_j^*$
is positive definite, and $d_j^{\rm gain}\geq1$.
The original polynomial remainder map in these orthonormal source
coordinates is exactly $C_sA_j$. This follows by applying it to
each displayed basis vector, including all high degrees.

For $x\in\mathbb C^q$, the unique minimum norm source lift of $C_sx$
at degree $q+j-1$ has coefficient vector
$A_j^*K_j^{-1}x$. Indeed its remainder coordinates are $x$,
and it is orthogonal to $\ker A_j$. Every other lift differs by
an element of that kernel, so its squared norm is
$x^*K_j^{-1}x$ plus the squared norm of that difference. Hence the
original quotient Gram is
\[
 G_{q+j-1}=C_s^{-*}K_j^{-1}C_s^{-1}.
 \tag{GMB4}
\]
This proves the actual quotient map and its metric, not only its
determinant. The source kernel has its full dimension $j$.

\subsection{The monic relation, its original phase, and its row}
Put $D=q+j$. In degree at most $D$, with the last coordinate
corresponding to $u_D$, form
\[
 z_j=\begin{pmatrix}-A_j^*K_j^{-1}v_j\\1\end{pmatrix},\qquad
 \mathfrak d_{D-1}(S)
   =i^D\sqrt{h_D^{(s)}}\left(
       u_D(S)-\sum_{a=0}^{D-1}(A_j^*K_j^{-1}v_j)_a u_a(S)
                         \right).
 \tag{GMB5}
\]
The first vector has zero remainder because
$-A_jA_j^*K_j^{-1}v_j+v_j=0$. Thus $\mathfrak d_{D-1}$ is
divisible by the original $\chi$. Its leading coefficient in $S$
is exactly one by GMB1. Moreover $z_j$ is orthogonal to every
old relation: its old coordinates lie in $\operatorname{im}A_j^*$.
It follows that $\mathfrak d_{D-1}=\chi\,p_j^{\chi,s}$, where
$p_j^{\chi,s}$ is the original monic polynomial orthogonal to lower
degrees in the measure $|\chi(c+iy)|^2dm_{0,s}(y)$. To verify
uniqueness without any assumption on coordinates, the difference
of two such monic relations is an old relation, and its inner
product with that difference is zero. Positivity makes it zero.
Consequently GMB5 is exactly the relation line used in BRU and
HMR32, with its prescribed monic orientation.

Let $a_{D-1}=\|\mathfrak d_{D-1}\|^2$ and let
$\alpha_{D-1}$ be the row with entries
$\langle\mathfrak d_{D-1},S^a\rangle$, $0\leq a<q$.
GMB3--5 give the complete formulas
\[
\begin{gathered}
 a_{D-1}=h_D^{(s)}d_j^{\rm gain},\\
 \alpha_{D-1}
   =-i^{-D}\sqrt{h_D^{(s)}}\,v_j^*K_j^{-1}C_s^{-1}.
\end{gathered}\tag{GMB6}
\]
For the first equality, the squared norm of $z_j$ is
$1+v_j^*K_j^{-1}A_jA_j^*K_j^{-1}v_j=1+t_j$.
For the second, the degree $D$ orthogonal polynomial has zero
inner product with every $S^a$, $a<q$. A polynomial of degree
less than $q$ has its low orthonormal coordinates $C_s^{-1}[P]$,
so the old lift term in GMB5 has inner product
$v_j^*K_j^{-1}C_s^{-1}[P]$ with it. Conjugating the scalar
$i^D\sqrt{h_D^{(s)}}$ proves precisely the phase and minus sign
in GMB6.

Direct multiplication proves
\[
 (K_j+v_jv_j^*)^{-1}
 =K_j^{-1}-\frac{K_j^{-1}v_jv_j^*K_j^{-1}}{d_j^{\rm gain}}.
\]
Using GMB4 and GMB6 therefore gives the original BRU update,
including its complete coefficient row:
\[
 G_{D-1}=G_D+\frac{\alpha_{D-1}^*\alpha_{D-1}}{a_{D-1}},
 \qquad
 \frac{\det G_{D-1}}{\det G_D}=d_j^{\rm gain}.
 \tag{GMB7}
\]
The determinant equality follows either by taking determinants
in GMB4, or by conjugating $K_j+v_jv_j^*$ by $K_j^{-1/2}$:
the resulting rank-one update has eigenvalue $1+t_j$ in its
one possible nontrivial direction and eigenvalue one elsewhere.
This also proves it when $v_j=0$.

\subsection{The unchanged observation and both mixed contractions}
The observation in these explicit source coordinates is
\[
 R_s=\Lambda C_s:\mathbb C^q\twoheadrightarrow B_{\rm obs},
 \quad M_j^{\rm obs}=R_sK_jR_s^*,\quad w_j=R_sv_j,\quad
 \xi_j=w_j^*(M_j^{\rm obs})^{-1}w_j.
 \tag{GMB8}
\]
The notation $R_s$ here names a matrix of the fixed observation,
not a new observation. Every Taylor unit and primary coefficient
already in $\Lambda$ is still in this product. For zero boundary
rank the matrices are empty and $\xi_j=0$. Otherwise
$M_j^{\rm obs}$ is positive definite by surjectivity. Define
\[
 r_j^{\rm ker}=v_j-K_jR_s^*(M_j^{\rm obs})^{-1}w_j.
 \tag{GMB9}
\]
Then $R_sr_j^{\rm ker}=0$. Its two summands give the exact
$K_j^{-1}$-orthogonal decomposition
\[
 t_j=\xi_j+
       (r_j^{\rm ker})^*K_j^{-1}r_j^{\rm ker},\qquad
 0\leq\xi_j\leq t_j.
 \tag{GMB10}
\]
Indeed the cross inner product is
$(r_j^{\rm ker})^*R_s^*(M_j^{\rm obs})^{-1}w_j=0$, and the
squared norm of the second summand in GMB9 is $\xi_j$.
The map $C_s$ restricts to a bijection
$\ker R_s\to\ker\Lambda$, since
$\Lambda C_sx=R_sx$. This supplies the exact map of the retained
kernel, including its original coefficient embedding.

Let $I:\ker\Lambda\hookrightarrow E_\chi$ be that original
embedding. At the new degree $D$, retain the original forms and
minimum section
\[
\begin{gathered}
 H_D^{\rm ker}=I^*G_DI,\qquad
 Q_D=(\Lambda G_D^{-1}\Lambda^*)^{-1},\qquad
 S_D=G_D^{-1}\Lambda^*Q_D,\\
 \beta_{D-1}=(\alpha_{D-1}I)(H_D^{\rm ker})^{-1}
                                  (\alpha_{D-1}I)^*,\qquad
 \nu_{D-1}=(\alpha_{D-1}S_D)Q_D^{-1}
                                  (\alpha_{D-1}S_D)^*.
\end{gathered}\tag{GMB11}
\]
All inverses on a zero dimensional space have their unique empty
meaning. The section obeys $\Lambda S_D=I$ and $I^*G_DS_D=0$.
The map $[I,S_D]$ is a bijection: the inverse sends $z$ to its
kernel component $z-S_D\Lambda z$ and its boundary component
$\Lambda z$. In this frame the metric has diagonal blocks
$H_D^{\rm ker}$ and $Q_D$. Inverting the congruence proves
\[
 G_D^{-1}=I(H_D^{\rm ker})^{-1}I^*
                       +S_DQ_D^{-1}S_D^*.
 \tag{GMB12}
\]
This proves all decomposition identities used below while
retaining the original source, kernel, boundary and section maps.

Set $d=1+t_j$, $h=h_D^{(s)}$, and
$N=M_j^{\rm obs}+w_jw_j^*$. Equations GMB4 and GMB8 imply
$Q_D=N^{-1}$ and
$S_D=C_s(K_j+v_jv_j^*)R_s^*N^{-1}$. Direct multiplication yields
$N^{-1}w_j=(M_j^{\rm obs})^{-1}w_j/(1+\xi_j)$.
GMB6 now gives
\[
 \alpha_{D-1}S_D
 =-i^{-D}\sqrt h\,d\,w_j^*N^{-1},\qquad
 \nu_{D-1}=h d^2\frac{\xi_j}{1+\xi_j}.
\]
On the other hand GMB6 and GMB12 give
\[
 \beta_{D-1}+\nu_{D-1}
   =\alpha_{D-1}G_D^{-1}\alpha_{D-1}^*
   =h(t_j+t_j^2)=h d t_j.
\]
Subtracting the displayed expression for $\nu_{D-1}$, with no
sign assumption, proves the exact full mixed row formula
\[
\boxed{\begin{gathered}
 a_{D-1}=h d,\qquad
 \beta_{D-1}=h d\frac{t_j-\xi_j}{1+\xi_j},\qquad
 \nu_{D-1}=h d^2\frac{\xi_j}{1+\xi_j},\\
 1+\frac{\beta_{D-1}}{a_{D-1}}
      =\frac{1+t_j}{1+\xi_j},\qquad
 1+\frac{\nu_{D-1}}{a_{D-1}+\beta_{D-1}}=1+\xi_j.
\end{gathered}}\tag{GMB13}
\]
Both factors are at least one by GMB10, and their product is
$1+t_j$, the original complete increment in GMB7. These are the
individual BRU kernel and boundary factors, in their original
order. In particular the kernel factor has the denominator
$1+\xi_j$; it is not replaced by $1+t_j-\xi_j$.

\subsection{The four original endpoints and their finite calculation}
For a sequence $f_j$, put
\[
 \mathcal W_q(f)=f_0+2\sum_{j=1}^{q-1}f_j+f_q.
 \tag{GMB14}
\]
The original signed quotient volume is
$\log\det G_{q-1}+\log\det G_q
 -\log\det G_{2q-1}-\log\det G_{2q}$.
Telescoping GMB7 at these four precise degrees gives
\[
\begin{gathered}
 \mathcal B_k^{(s)}=\mathcal W_q\bigl(\log(1+t_j)\bigr),\\
 F_{K}^{(s)}=\mathcal W_q\left(
                    \log\frac{1+t_j}{1+\xi_j}\right),\qquad
 F_{B}^{(s)}=\mathcal W_q\bigl(\log(1+\xi_j)\bigr),\\
 F_{K}^{(s)}+F_{B}^{(s)}=\mathcal B_k^{(s)}.
\end{gathered}\tag{GMB15}
\]
To check the indices, the relation added at this step has degree
$D=q+j$ and is the BRU row with index $D-1$. Thus $j=0,q$
are exactly the two weight-one rows $q-1,2q-1$; all intervening
rows have weight two. GMB13 proves the individual factors, so
GMB15 is the full original mixed decomposition at $s=1$ or $s=k$.
It does not assign an arithmetic source to either Gamma order.
The established arithmetic comparison retains its own exact
transition scalar and its original maps.

For completeness this is a finite calculation from the original
polynomial, rather than a definition by an unspecified projection.
Let $T_\chi$ be multiplication by $S$ on the fixed remainder frame
and set $Y_\chi=(T_\chi-cI)/i$. Starting with
$r_0=[1]$, $r_1=Y_\chi[1]$, compute
\[
 r_{d+1}=Y_\chi r_d-d(b_s+d-1)r_{d-1},\qquad
 C_s=\left[\frac{r_0}{\sqrt{h_0^{(s)}}},\ldots,
                 \frac{r_{q-1}}{\sqrt{h_{q-1}^{(s)}}}\right],
 \qquad v_j=\frac{C_s^{-1}r_{q+j}}{\sqrt{h_{q+j}^{(s)}}}.
 \tag{GMB16}
\]
Polynomial division and GMB1 prove by induction that $r_d=[p_d(Y)]$.
Every entry of $T_\chi$ is a literal coefficient of the original
$\chi$, including its primary multiplicities. The matrices
$K_j$, $M_j^{\rm obs}$ and the scalars in GMB13--15 are then obtained
by the displayed finite operations. Their inverses exist by the
proved positivity, including the stated empty cases. This closes
the exact morphism from the intrinsic exterior mass calculation,
through the original polynomial remainders, to both original
mixed contractions. It introduces no new source order or factor.

% END INLINE SOURCE 28

\clearpage

% BEGIN INLINE SOURCE 29; SHA256 517daf9230b256ed0fad88e5104ba9fd8f8d51e793998fea955c1a34bacf480d
\section{The boundary and its retained kernel under each original relation}
\label{GraphV3:reader-v3-BRU:sec:bru-original-relations}

This calculation continues the incoming boundary observation A1 and the
boundary/kernel metric A4--A18. It calculates those metrics and every
cross-degree overlap from the original relation rows. All maps below are
on the actual packet, before taking a completion. In particular the
boundary observation does not remove its kernel.

\subsection{The packet, observation, source and finite relation rows}

Retain an actual hypothetical complete-order quartet
\[
 \rho=\tfrac12+\delta+i\gamma,\quad 0<\delta<\tfrac12,\quad\gamma>2,
 \quad h(s)=\prod_{r\in\{\rho,\bar\rho,1-\rho,1-\bar\rho\}}(s-r)^m,
 \quad n=4m,
 \tag{BRU1}
\]
where $m\geq1$ is its complete zero order in $g=2\xi$.
Thus $v_h=g/h$ is the original entire amplitude and
$U=M_{j_hv_h}$ is its invertible multiplication matrix on
$E_h=\mathbb{C}[s]/(h)$. No value of this unit is assigned.
For $k\geq1$ the original cyclic coefficient space and action are
\[
 \begin{gathered}
 \ell_k=1+k(m-1),\qquad q=(k+1)^2\ell_k,\qquad E=\mathbb{C}[S]/(\chi),\\
 \chi(S)=\prod_{a,b=0}^k
 (S-k/2-(2a-k)\delta-i(2b-k)\gamma)^{\ell_k},\qquad
 A=M_S,\\
 \iota:E\longrightarrow E_h^{\otimes k},\qquad
 [P]\longmapsto[P(s_1+\cdots+s_k)],\qquad
 \eta=U^{\otimes k}\iota .
 \end{gathered}\tag{BRU2}
\]
The original cyclic inclusion is injective: on each ordered primary
tensor the sum of the nilpotents has nilpotency index
$1+k(m-1)$, since its power $k(m-1)$ has nonzero coefficient
$(k(m-1))!/((m-1)!)^k$ on the product of the top powers. Its next
power is zero by total degree. All sums in BRU2 are distinct for
different $(a,b)$ because $\delta,\gamma>0$. Each occurs among the
ordered factors, so their product with exactly these exponents is
the minimal polynomial of the tensor sum. This proves the stated
kernel of polynomial evaluation and injectivity of $\iota$.

Fix $u\ne0$ on the original logarithm branch and $t=0$. Let
$\Pi$ be the original period matrix with rays
$\Gamma_j=-\ell_0+\ell_j$, angles
$(\arg u+\pi+2\pi j)/(n+1)$, and entries
$\int_{\Gamma_j}e^{\Phi_h(s)/u}s^b\,ds$, where $\Phi_h'=h$ and
$\Phi_h(0)=0$. Its full determinant and convergence are proved in
PDE1--13 and IPM1--7, included in the source package. Let $C$ have
rows describing $\Gamma_j\mapsto\Gamma_{j+1}-\Gamma_1$ for $j<n$
and $\Gamma_n\mapsto-\Gamma_1$. Set
\[
 P_k=\frac1{n+1}\sum_{a=0}^{n}(C^{-a})^{\otimes k},\qquad
 L=P_k\Pi^{\otimes k}\eta:E\longrightarrow(\mathbb{C}^n)^{\otimes k}.
 \tag{BRU3}
\]
For the algebra used here, $C^{n+1}=I$ follows by applying the
cyclic permutation to the reduced endpoint basis $e_j-e_0$.
Multiplying the finite sum therefore gives $P_k^2=P_k$.
We use the literal coefficient map BRU3, without evaluating a
formal stationary series or identifying its action with Frobenius.
Choose a basis inclusion $\jmath:B=\operatorname{im}L\hookrightarrow
(\mathbb{C}^n)^{\otimes k}$, and define the surjection
$\Lambda:E\to B$ by $\jmath\Lambda=L$. Its kernel is $K$.
Fix a basis inclusion $I:K\hookrightarrow E$, and a coefficient
section $S_*$ with $\Lambda S_*=I_B$. Gaussian elimination on the
specified finite matrix constructs these maps, including at a
rank drop by using its actual image there. Keep these choices fixed
as the source metric varies. Write $r=\dim B$ and $d_K=q-r$.
Zero-dimensional factors use their unique empty inverse and determinant one.


% END INLINE SOURCE 29

\clearpage

% BEGIN INLINE SOURCE 30; SHA256 347a84b27a1b7cc9fee7f5c4a155400c1e3030a7a1f69c84400438e9fb21d772
\section{The actual entire theta amplitude and the finite period observation}
\label{sec:ipm-entire-source}

This calculation uses the actual packet
\(h(s)=\prod_\rho(s-\rho)^{m_\rho}\), with all complete orders,
\(d_h=\deg h\ge1\), \(\Phi_h'=h\), \(\Phi_h(0)=0\),
\(g=2\xi\), \(v_h=g/h\), and
\(\upsilon_h=j_hv_h\in E_h^\times\), where
\(E_h=\mathbb C[s]/(h)\). No root outside this actual packet is
introduced. The analytic statements have these data as their domain;
they do not assert that an off-critical packet has been found.
The finite-source map is the original injective map
\(\mathcal T_hP=P(D)F_h\), with Mellin transform \(v_hP\).
Thus all maps below on \(\mathcal T_h\mathcal P_N\) are defined
unambiguously through that injective source parametrization.

\subsection{An explicit growth bound proves the global integration domain}

Write
\[
 \psi(x)=\sum_{n\ge1}e^{-\pi n^2x},\qquad
 C_\theta=(1-e^{-3\pi})^{-1},\qquad
 \beta_R=(R+1)/2.
 \tag{IPM.1}
\]
The original theta representation, also recorded in DLMF 25.5.13--14,
gives the entire identity
\[
 g(s)=1+s(s-1)\int_1^\infty
       \bigl(x^{s/2}+x^{(1-s)/2}\bigr)\psi(x)\,\frac{dx}{x}.
 \tag{IPM.2}
\]
For completeness, the identity follows initially for \(\Re s>1\)
by integrating each Gaussian:
\(\int_0^\infty\psi(x)x^{s/2}dx/x
=\pi^{-s/2}\Gamma(s/2)\zeta(s)\).
Split the integral at one and use the original Poisson identity
\(\psi(x)=x^{-1/2}\psi(1/x)+(x^{-1/2}-1)/2\).
The elementary integral on \((0,1)\) is
\(1/(s-1)-1/s=1/(s(s-1))\). Multiplication by \(s(s-1)\)
proves (IPM.2) in that half-plane. Its right side is entire, since its
integral and every coefficient derivative converge uniformly on compact
sets by Gaussian decay. Analytic continuation proves the identity
everywhere, including \(s=0,1\).

Since \(n^2\ge1+3(n-1)\), for \(x\ge1\)
one has \(\psi(x)\le C_\theta e^{-\pi x}\). Consequently
\[
 \begin{gathered}
 |g(s)|\le\mathcal X(R):=
  1+2C_\theta R(R+1)\pi^{-\beta_R}\Gamma(\beta_R,\pi)\\
 \le 1+2C_\theta R(R+1)\pi^{-\beta_R}\Gamma(\beta_R),
 \qquad |s|\le R,\ R\ge1.
 \end{gathered}
 \tag{IPM.3}
\]
Indeed both powers of \(x\) in the integral have absolute value at most
\(x^{(R+1)/2}\), and \(|s(s-1)|\le R(R+1)\).
The upper incomplete gamma is the literal integral after \(z=\pi x\).
An elementary bound sufficient here is
\(\Gamma(\beta)\le2(2\beta)^\beta\) for \(\beta\ge1\).
To prove it, write its integrand as
\((t^{\beta-1}e^{-t/2})e^{-t/2}\), bound the first factor by
its maximum \([2(\beta-1)/e]^{\beta-1}\), and integrate the second.
At \(\beta=1\) use the value one for the zeroth power. This maximum
is at most \((2\beta)^\beta\). In particular
\(\log\mathcal X(R)=O(R\log(R+2))\); the explicit function
\(\mathcal X\), rather than a substituted constant, is retained in
all tail integrals.

Let \(R_h=\max_\rho|\rho|\). For \(|s|=r\ge\max(1,2R_h)\),
\[
 |h(s)|\ge(r/2)^{d_h},\qquad
 |v_h(s)|\le 2^{d_h}r^{-d_h}\mathcal X(r).
 \tag{IPM.4}
\]
The complete-divisor hypothesis makes \(v_h\) entire, so it is bounded
on the remaining compact disc. Cauchy's integral formula on radius-one
circles gives the same \(\exp(O(r\log(r+2)))\) growth for each fixed
derivative. The formula retains both the degree and every original root
in its compact-disc and exterior bounds.

Fix \(u\ne0\), \(t\in\mathbb C\), a branch of \(\arg u\), and
the original BC rays and oriented contours
\[
 \theta_j=\frac{\arg u+\pi+2\pi j}{d_h+1},\quad
 \ell_j(r)=re^{i\theta_j},\quad
 \Gamma_j=-\ell_0+\ell_j\quad(1\le j\le d_h),\quad
 f_t=\Phi_h-ts.
 \tag{IPM.5}
\]
Their leading real exponent is
\(-r^{d_h+1}/((d_h+1)|u|)\). If
\(f_t(s)=s^{d_h+1}/(d_h+1)+\sum_{\nu=0}^{d_h}c_\nu s^\nu\),
then, for \(r\ge1\), its full real exponent is at most
\[
 -\frac{r^{d_h+1}}{(d_h+1)|u|}
       +\frac{\sum_{\nu=0}^{d_h}|c_\nu|r^\nu}{|u|}.
 \tag{IPM.6}
\]
Combining (IPM.3)--(IPM.6) proves absolute convergence of
\[
 \mathcal I_{u,t}(a)=
  \left(\int_{\Gamma_j}e^{f_t(s)/u}a(s)\,ds\right)_{j=1}^{d_h}
 \tag{IPM.7}
\]
for \(a=v_hP\), every polynomial \(P\), every derivative used below,
and every entire quotient by \(h\) constructed below. On a compact
parameter set with \(u\ne0\), the same domination proves coefficient
differentiation and the needed contour deformations inside the fixed
decaying sectors. This is a proof of the domain for the actual theta
amplitude. It makes no assertion for arbitrary undeclared rays.


% END INLINE SOURCE 30

\clearpage

% BEGIN INLINE SOURCE 31; SHA256 4df10819db4933332fc6d72992b26b1c9a9f0da11e1d8c813c1c0611e31a9f9c
% Exact scalar dependency used in PDE9--11 and IPM.20--21.
\section{The positive Gamma grid product in the original period determinant}
\label{sec:PGCC}

The full original PDE determinant proof reduces its constant to
\[
 \prod_{r=1}^{d-1}\Gamma(r/d)
       =\frac{(2\pi)^{(d-1)/2}}{\sqrt d},\qquad d=n+1\geq2.
 \tag{PGCC1}
\]
Here is a complete proof of this particular scalar identity from the
positive defining Gamma integral, retaining its positive square root.
For a real $0<a<1$, Tonelli's theorem applies to the nonnegative
integrand defining the product of two Gamma integrals. Set
$x=st/(1+t)$ and $y=s/(1+t)$, with $s,t>0$. This is a bijective
change of variables on the positive quadrant with Jacobian
$s/(1+t)^2$, and direct multiplication gives
\[
 \begin{split}
 \Gamma(a)\Gamma(1-a)
 &=\int_0^\infty\int_0^\infty
        e^{-x-y}x^{a-1}y^{-a}\,dx\,dy\\
 &=\int_0^\infty e^{-s}\,ds
       \int_0^\infty\frac{t^{a-1}}{1+t}\,dt
  =\int_0^\infty\frac{t^{a-1}}{1+t}\,dt.
 \end{split}\tag{PGCC2}
\]
The last integral converges at zero because $a>0$ and at infinity
because $a<1$.

To evaluate it with its exact sign, use
$z^{a-1}=\exp((a-1)\operatorname{Log}z)$ with $0<\arg z<2\pi$.
The positively oriented keyhole contour about the positive real axis
has an upper segment from $\varepsilon$ to $R$, an outer circle
counterclockwise, a lower segment from $R$ to $\varepsilon$, and an
inner circle clockwise. Take $0<\varepsilon<1<R$ and let the banks
approach the cut. The only pole of $z^{a-1}/(1+z)$ inside is $z=-1$,
with residue $e^{\pi i(a-1)}$. The two straight segments sum to
\[
 (1-e^{2\pi ia})\int_\varepsilon^R
                       \frac{t^{a-1}}{1+t}\,dt.
 \tag{PGCC3}
\]
The outer circular contribution has absolute value at most
$2\pi R^a/(R-1)$, which tends to zero as $R\to\infty$.
The inner contribution has absolute value at most
$2\pi\varepsilon^a/(1-\varepsilon)$, which tends to zero as
$\varepsilon\to0$. The residue theorem and these bounds give
\[
 \begin{split}
 (1-e^{2\pi ia})\int_0^\infty\frac{t^{a-1}}{1+t}\,dt
 &=2\pi i e^{\pi i(a-1)},\\
 \Gamma(a)\Gamma(1-a)&=\frac{\pi}{\sin(\pi a)}.
 \end{split}\tag{PGCC4}
\]
In the second line the quotient sign follows explicitly from
$1-e^{2\pi ia}=-2i e^{\pi ia}\sin(\pi a)$ and
$e^{\pi i(a-1)}=-e^{\pi ia}$.

Now put $\zeta=e^{2\pi i/d}$. Factoring $X^d-1$ over its distinct
roots and evaluating the quotient by $X-1$ at $X=1$ proves
\[
 d=\prod_{r=1}^{d-1}(1-\zeta^r),\qquad
 d=\prod_{r=1}^{d-1}|1-\zeta^r|
   =2^{d-1}\prod_{r=1}^{d-1}\sin(\pi r/d).
 \tag{PGCC5}
\]
The sines in this product are all positive. Apply PGCC4 at
$a=r/d$ and multiply over $1\leq r<d$. The involution
$r\mapsto d-r$ permutes the same factors, so
\[
 \left(\prod_{r=1}^{d-1}\Gamma(r/d)\right)^2
 =\prod_{r=1}^{d-1}\frac{\pi}{\sin(\pi r/d)}
 =\frac{(2\pi)^{d-1}}d.
 \tag{PGCC6}
\]
Each Gamma integral is strictly positive. Taking the positive square
root proves PGCC1, including the case $d=2$.

This evaluates the precise scalar used in PDE9--11 without replacing
the original Fourier determinant or its phase. In fact the polynomial
period isomorphism already follows before this evaluation: PDE9 gives
nonzero positive Gamma factors and fixed nonzero powers of $u$, while
PDE10 gives the nonzero Fourier Vandermonde determinant. PDE6--8 then
propagate that nonzero determinant through the full coefficient family.
PGCC1 additionally supplies its displayed value $(2\pi)^{n/2}$,
with all original powers and phases calculated in PDE9--11 unchanged.

% END INLINE SOURCE 31


% BEGIN INLINE SOURCE 32; SHA256 bc8cc467be1b3858772131152cb2823968f886d9aae34baa845e71483cec97bb
\section{The full polynomial period deformation and its entire-amplitude forcing}
\label{GraphV3:sec:pde-full-deformation}

Retain the original nonempty complete packet $h(s)$, its degree $n\geq1$,
its primitive $\Phi_h(0)=0$, its arithmetic coefficient space
$E_h=\mathbb{C}[s]/(h)$ in the increasing frame $1,s,\ldots,s^{n-1}$,
and the actual entire amplitude $v_h=2\xi/h$. All zero orders in $h$
are complete, so this last quotient is entire and its Taylor unit
$\upsilon_h=j_h(v_h)$ is invertible in $E_h$. Write
\[
 h(s)=s^n+\sum_{r=0}^{n-1}h_rs^r,\qquad
 f_t(s)=\Phi_h(s)-ts,\qquad
 A_h(t)=A_h+tR,\quad R=e_0e_{n-1}^{\mathsf T}.
 \tag{PDE1}
\]
Here $A_h$ is the original arithmetic multiplication by $s$ on $E_h$;
$A_h(t)$ is its displayed companion deformation on the same coefficient
frame. The parameter $t$ does not replace the original marked arithmetic
quotient by a different quotient. The matrix $R$ includes the case
$n=1$, where it is the identity.

Fix $u\ne0$ on a specified branch of $\log u$. Set $d=n+1$ and retain
the actual rays and their orientations
\[
 \theta_j=\frac{\arg u+\pi+2\pi j}{d},\quad
 \ell_j(r)=re^{i\theta_j}\ (r\geq0),\quad
 \Gamma_j=-\ell_0+\ell_j\quad(1\leq j\leq n).
 \tag{PDE2}
\]
For an admitted amplitude $a$, let $\mathcal I_t(a)$ be the column with
entries $\int_{\Gamma_j}e^{f_t/u}a(s)\,ds$. Every amplitude used below is
a polynomial times a derivative of the fixed $v_h$, or a polynomial.
IPM1--6 proves their convergence on these rays with the original theta
bound; the explicit derivative estimate below also proves the required
first-derivative convergence directly. Compact sets of parameters have
a common dominating function with negative leading term
$-r^d/(d|u|)$ and lower polynomial terms. Thus these integrals and their
displayed parameter derivatives are entire in $t$. The fixed rays have
no $t$-dependent endpoint or orientation.

\subsection{The exact determinant on the whole polynomial family}

For this calculation only, allow all lower phase coefficients to vary:
\[
 f_{\mathbf c}(s)=\frac{s^{n+1}}{n+1}+\sum_{j=0}^n c_js^j,
 \quad h_{\mathbf c}=f_{\mathbf c}',\quad
 L_{\mathbf c}=u\partial_s+h_{\mathbf c},\quad
 \Pi(\mathbf c,u)_{j,b}
       =\int_{\Gamma_j}e^{f_{\mathbf c}/u}s^b\,ds.
 \tag{PDE3}
\]
The row index is $1\leq j\leq n$ and the column index $0\leq b<n$.
The original phase $f_t$ is the specified coefficient slice of this
family. No coefficient of that slice will be suppressed.

Ordinary monic division gives, for $0\leq j\leq n$ and $0\leq b<n$,
\[
 s^{b+j}=h_{\mathbf c}q_{b,j}+r_{b,j},\qquad
 \deg r_{b,j}<n,\quad\deg q_{b,j}\leq b+j-n.
 \tag{PDE4}
\]
A negative bound means the zero polynomial. In particular
$\deg q_{b,j}'\leq b-1<n$, so the exact differential reduction is
\[
 s^{b+j}=L_{\mathbf c}q_{b,j}
                    +(r_{b,j}-u q_{b,j}').
 \tag{PDE5}
\]
The leading term of $L_{\mathbf c}P$ has degree $n+\deg P$ and the same
leading coefficient as $P$. Repeated subtraction therefore proves that
the cokernel of $L_{\mathbf c}$ has the unique increasing frame of
degrees less than $n$, for every $\mathbf c$ and $u$. Formula PDE5 is
its literal remainder, not a specialization at $u=0$.

Let $B_j(\mathbf c,u)$ have these remainder columns, and let
$M_{\mathbf c}$ be multiplication by $s$ on
$\mathbb{C}[s]/(h_{\mathbf c})$. Ordinary remainders form
$M_{\mathbf c}^j$. The correction column $-u q_{b,j}'$ has row index
strictly less than $b$. It is consequently strictly triangular and has
zero trace. Thus, retaining its full nonzero matrix,
\[
 B_j=M_{\mathbf c}^{j}-u(\operatorname{coeff}q_{b,j}')_{b=0}^{n-1},
 \qquad \operatorname{Tr}B_j=\operatorname{Tr}M_{\mathbf c}^{j}.
 \tag{PDE6}
\]
The boundary integral of $e^{f_{\mathbf c}/u}q_{b,j}$ vanishes at
the two outer ends and cancels at the common origin. Differentiating
the original integrals and applying PDE5 gives
\[
 \partial_{c_j}\Pi=\frac1u\Pi B_j\quad(0\leq j\leq n).
 \tag{PDE7}
\]

Define the polynomial in the original coefficients
\[
 \mathcal S(\mathbf c)
  =\operatorname{Tr} f_{\mathbf c}(M_{\mathbf c})
  =\frac{\operatorname{Tr}M_{\mathbf c}^{n+1}}{n+1}
                +\sum_{j=0}^n c_j\operatorname{Tr}M_{\mathbf c}^j.
 \tag{PDE8}
\]
Its derivative with respect to $c_j$ is
$\operatorname{Tr}M_{\mathbf c}^j$. Here is a proof including repeated
roots. On the open set where $h_{\mathbf c}$ has distinct roots
$r_1,\ldots,r_n$, they admit local holomorphic enumerations and
$\mathcal S=\sum_i f_{\mathbf c}(r_i)$. Differentiation gives
$\partial_{c_j}\mathcal S=\sum_i r_i^j$, since
$f_{\mathbf c}'(r_i)=0$ makes every root-derivative term exactly zero.
Both sides of the identity are polynomials in $\mathbf c$; the open
set is dense because the derivative $s^n-1$ occurs in this family
and has distinct roots. Equality extends as a polynomial identity to
every coefficient value. Equivalently, the trace at a repeated root
is its full multiplicity times the scalar value, since powers of its
nilpotent part have zero diagonal. No multiplicity is removed.

Put $D=\det\Pi$. Multilinearity of the determinant in its columns
gives $\partial_{c_j}D=u^{-1}(\operatorname{Tr}B_j)D$ without any
prior assumption that $\Pi$ is invertible. PDE6 and PDE8 imply that
every derivative of $D\exp(-\mathcal S/u)$ is zero on the connected
coefficient space. Its value is determined at $\mathbf c=0$.

For completeness, write $\zeta=\exp(2\pi i/d)$ and
$\beta_b=(b+1)/d$. Direct integration on the original two rays gives
\[
 \Pi(0,u)_{j,b}
  =u^{(b+1)/d}d^{\beta_b-1}
          e^{i\pi\beta_b}\Gamma(\beta_b)
                       (\zeta^{j(b+1)}-1).
 \tag{PDE9}
\]
The power of $u$ uses the branch fixed in PDE2. To compute the remaining
determinant let $V_{r,b}=\zeta^{rb}$, $0\leq r,b<d$. Subtract row zero
from each other row and expand at its first column. The resulting
minor is exactly $(\zeta^{j(b+1)}-1)_{1\leq j\leq n,0\leq b<n}$,
with no row or column reversal. The Vandermonde formula gives
\[
 \det V=\prod_{0\leq p<q<d}(\zeta^q-\zeta^p),\qquad
 \zeta^q-\zeta^p
       =2i e^{\pi i(p+q)/d}\sin\frac{\pi(q-p)}d.
 \tag{PDE10}
\]
Every displayed sine is positive. The phase of this product is
$\pi d(d-1)/4+\pi(d-1)^2/2
=\pi(d-1)(3d-2)/4$.
Its first term comes from the $d(d-1)/2$ factors $i$, and its second
from $\sum_{p<q}(p+q)=d(d-1)^2/2$.
Also $V^*V=dI$ by the finite geometric sum, so $|\det V|=d^{d/2}$.
The retained gamma product is
$\prod_{b=0}^{n-1}\Gamma((b+1)/d)=(2\pi)^{n/2}/\sqrt d$,
the multiplication identity used in IPM21, also recorded in DLMF
5.5.7 at \url{https://dlmf.nist.gov/5.5.E7}.
The powers of $d$ from PDE9 sum to $-n/2$ and cancel this product's
$d^{-1/2}$ against $d^{d/2}$. Its additional phase is
$\pi\sum_b\beta_b=\pi n/2$. The total phase is therefore
$3\pi n(n+1)/4$. Consequently the exact determinant is
\[
 \boxed{\displaystyle
 \det\Pi(\mathbf c,u)
   =(2\pi)^{n/2}u^{n/2}
       \exp\left(\frac{3\pi i n(n+1)}4
                         +\frac{\mathcal S(\mathbf c)}u\right).}
 \tag{PDE11}
\]
In particular it never vanishes for $u\ne0$, including coefficients
with repeated critical points or repeated critical values. This proves
invertibility rather than assuming it to use a determinant equation.
For $n=1$ the formula is the downward Gaussian for $u>0$ on
$\arg u=0$, continued on the specified branch:
$-i\sqrt{2\pi u}\exp((c_0-c_1^2/2)/u)$.

\subsection{The original full quartet and the exact complex-parameter transport}

For the original full quartet $n=4m$, retain
$\rho=1/2+\delta+i\gamma$, $0<\delta<1/2$, $\gamma>2$, and
$h=\prod_{r\in\{\rho,\bar\rho,1-\rho,1-\bar\rho\}}(s-r)^m$.
Its reflection is $h(1-s)=h(s)$. With the original primitive constant
$A=\Phi_h(1/2)$, differentiation and evaluation at $s=1/2$ prove
\[
 \Phi_h(1-s)=2A-\Phi_h(s),\qquad
 f_t(1-s)=2A-t-f_t(s),\qquad
 \mathcal S(f_t)=n(A-t/2).
 \tag{PDE12}
\]
For the trace assertion, reflection pairs the roots of $h(s)-t$ with
their full multiplicities. Each pair has the displayed sum of phase
values. A root fixed by reflection has $s=1/2$ and phase value
$A-t/2$ itself. Summing over all $n$ roots, or taking the trace of the
corresponding algebra involution, proves the identity even at every
multiple-root parameter. Hence on this original quartet
\[
 \det\Pi_h(t,u)
  =(2\pi)^{n/2}u^{n/2}
       e^{3\pi i n(n+1)/4}\,e^{n(A-t/2)/u}.
 \tag{PDE13}
\]
Neither the original $A$ nor the $t$ term has been discarded.


% END INLINE SOURCE 32

\end{document}
